% GENERATED FILE. Edit canonical lessons, then run npm run generate:books.
% canonical-source-hash: fnv1a64:033e6f024cba0da8
% canonical-lessons: IT-C263-parlavo, IT-C263-era, IT-C263-faceva, IT-C263-dormivo, IT-C263-ripasso-mentre

\chapter{The Scene --- Ero, Era, C'era}
\label{ch:it-the-scene}
\hlchaptermodality{263}

\begin{chapteropening}
\textbf{By the end of this chapter:} \emph{I can paint the scene of a past story with the imperfetto, and drop the events into it with the passato prossimo.}

Complete IT-C263-ripasso-mentre: a scene in the imperfetto and an event that cuts into it, joined with mentre.
\end{chapteropening}

\section[parlavo, avevo, dormivo]{parlavo — the past that is still running}

\label{lesson:IT-C263-parlavo}



\begin{quote}
Say \textquotedblleft{}I spoke\textquotedblright{} the way chapter fifteen built it. (\emph{Ho parlato.})

That past has an edge: it happened, and it finished. Italian has a second past with no edge at all, and every story you tell from here on needs both of them.
\end{quote}

\begin{grammarlens}[title={Grammar Lens: one ending, three families}]
Take the infinitive, drop its last two letters, \textbf{-re}, and add \textbf{-vo} for \textquotedblleft{}I\textquotedblright{}:

\begin{quote}
parla-re, \textbf{parlavo} — I was speaking, I used to speak ave-re, \textbf{avevo} — I had, I used to have dormi-re, \textbf{dormivo} — I was sleeping
\end{quote}

The vowel in front of the \textbf{-v-} is the family's own: \emph{-a-}, \emph{-e-}, \emph{-i-}. Nothing else moves. The other persons keep that \textbf{-v-} and change only what comes after it: \textbf{parlavo}, \textbf{parlavi}, \textbf{parlava}, \textbf{parlavamo}, \textbf{parlavate}, \textbf{parlavano} — I, you, he or she, we, you all, they.

This is the \textbf{imperfetto}, and it is the most regular tense Italian owns: one recipe for almost every verb. The few that refuse it begin in the next lesson.

What it means is \textquotedblleft{}was ‑ing\textquotedblright{}, or \textquotedblleft{}used to\textquotedblright{}. \emph{Parlavo} is speaking with no start and no finish in view, which is exactly what a story needs for its background.
\end{grammarlens}

\begin{cousinweb}
Latin already had this past, and it marked it with a \textbf{b}: Latin \emph{amābam}, \textquotedblleft{}I was loving\textquotedblright{}, and \emph{habēbat}, \textquotedblleft{}he was having\textquotedblright{}. Between two vowels that \emph{b} softened to a \emph{v}, the same softening that turned Latin \emph{habēre} into \textbf{avere}. So \emph{habēbat} comes down to \textbf{aveva} sound by sound.

One piece is younger. Old Italian said \emph{io parlava}, with the same \emph{-a} as \emph{lui parlava}. The \textbf{-o} of \emph{parlavo} was borrowed later from the present, \emph{parlo}, so that \textquotedblleft{}I\textquotedblright{} had an ending of its own. Neither EtymDB nor Lewis \& Short records endings this late; that step is the usual explanation in histories of Italian, and it is given here as one.
\end{cousinweb}

\subsection*{Your turn}
Say these aloud:

\begin{itemize}
\raggedright
  \item \emph{parlavo, avevo, dormivo} — one \emph{-v-}, three vowels in front of it
  \item \emph{parlavo, parlavi, parlava, parlavamo, parlavano}
  \item \emph{lavoravo} — chapter five's verb, still running in the past
  \item the edge, then no edge — \emph{ho parlato}, then \emph{parlavo}
\end{itemize}

\begin{tcolorbox}[breakable,colback=teal!4,colframe=teal!35!black,title={Before you move on}]
How do you build \textquotedblleft{}I was sleeping\textquotedblright{}? (Drop the \emph{-re} of \emph{dormire} and add \emph{-vo}: \emph{\textbf{dormivo}}.) Which Latin letter became the \emph{-v-}? (The \emph{\textbf{b}} of \emph{amābam}, softened as in \emph{habēre} and \textbf{avere}.) Which past has an edge, and which has none? (\emph{\textbf{Ho parlato}} has one; \emph{\textbf{parlavo}} has none.)
\end{tcolorbox}
\section[era, c'era]{era, c'era — the two words a scene is made of}

\label{lesson:IT-C263-era}



\begin{quote}
Build \textquotedblleft{}I was speaking\textquotedblright{} from \emph{parlare}. (\emph{Parlavo.}) Now say \textquotedblleft{}there is\textquotedblright{}. (\emph{C'è.})

Hold on to both. This lesson joins them.
\end{quote}

\subsection*{You'll want to know: era}
One verb refuses last lesson's recipe, and it is the one you use most: \textbf{essere}. Its imperfetto is short, and it is old:

\begin{quote}
\textbf{ero} — I was · \textbf{eri} — you were · \textbf{era} — he, she or it was \textbf{eravamo} — we were · \textbf{eravate} — you all were · \textbf{erano} — they were
\end{quote}

\textbf{Era} is the form you will hear most, because a scene is mostly about things and places rather than about you:

\begin{quote}
\emph{La stazione era piccola.} — The station was small. \emph{Ero stanco, ma ero contento.} — I was tired, but I was happy.
\end{quote}

\subsection*{You'll want to know: c'era}
You have \textbf{c'è} and \textbf{ci sono}, \textquotedblleft{}there is\textquotedblright{} and \textquotedblleft{}there are\textquotedblright{}. Move them into the imperfetto and you have the way a scene opens:

\begin{quote}
\textbf{c'era} — there was · \textbf{c'erano} — there were \emph{C'era il sole.} — The sun was out. (Literally, \textquotedblleft{}there was the sun\textquotedblright{}.) \emph{C'erano molte persone.} — There were a lot of people.
\end{quote}

Same \textbf{ci}, same apostrophe where it meets a vowel. Only the verb has moved into the past.

\begin{cousinweb}
\textbf{era} $\leftarrow$ Latin \emph{erat}, \textquotedblleft{}he was, she was\textquotedblright{} — inherited, with only the final \emph{t} worn away. Latin's past of \textquotedblleft{}to be\textquotedblright{} was irregular already, \emph{eram, erās, erat}, and it never had the \emph{b} of \emph{amābam}, so it never took the \emph{-v-} either. That is why this one form looks nothing like \emph{parlava}. (Lewis \& Short, under \emph{sum}.)

The \emph{ero} for \textquotedblleft{}I\textquotedblright{} got its \emph{-o} late, on the same usual account as \emph{parlavo}: old Italian said \emph{io era}. EtymDB's only row for an Italian \emph{era} is the noun \textquotedblleft{}an era, an age\textquotedblright{}, a different word altogether, from Latin \emph{aera}.
\end{cousinweb}

\subsection*{Your turn}
Say these aloud:

\begin{itemize}
\raggedright
  \item \emph{ero, eri, era, eravamo, erano}
  \item \emph{C'era il sole.}
  \item \emph{Ero stanco, ma ero contento.}
  \item from now to then — \emph{c'è} becomes \emph{c'era}, \emph{ci sono} becomes \emph{c'erano}
\end{itemize}

\begin{tcolorbox}[breakable,colback=teal!4,colframe=teal!35!black,title={Before you move on}]
Which Latin word is \emph{era}? (\emph{\textbf{Erat}}, \textquotedblleft{}he was\textquotedblright{} — only the \emph{t} lost.) Why does it not end in \emph{-va} like \emph{parlava}? (Latin's own form never had the \emph{b}.) Say \textquotedblleft{}there were a lot of people\textquotedblright{}. (\emph{\textbf{C'erano molte persone.}})
\end{tcolorbox}
\section[faceva freddo, faceva caldo]{faceva freddo — the weather, made in the past}

\label{lesson:IT-C263-faceva}



\begin{quote}
Say \textquotedblleft{}the sun was out\textquotedblright{} the way last lesson built it. (\emph{C'era il sole.}) Then \textquotedblleft{}I was tired\textquotedblright{}. (\emph{Ero stanco}, or \emph{ero stanca}.)
\end{quote}

\subsection*{You'll want to know: faceva freddo}
English says the weather \textbf{was} cold. Italian says it \textbf{made} cold:

\begin{quote}
\textbf{Faceva freddo.} — It was cold. \textbf{Faceva caldo.} — It was hot.
\end{quote}

\emph{Faceva} is the imperfetto of a verb meaning \textquotedblleft{}to make, to do\textquotedblright{}, and for heat and cold that is the verb Italian reaches for. It needs no subject: nobody is doing the making.

Keep \textbf{era} for things. \emph{Il caffè era freddo} — the coffee was cold — is about the coffee. \emph{Faceva freddo} is about the whole day.

And rain takes this chapter's plain recipe: \emph{piovere} gives \textbf{pioveva}, \textquotedblleft{}it was raining\textquotedblright{}. Put three of them together and you have a whole morning before anything has happened in it:

\begin{quote}
\emph{Faceva freddo e pioveva. Non c'era il sole.}
\end{quote}

\begin{cousinweb}
\textbf{faceva} $\leftarrow$ Latin \emph{faciēbat}, \textquotedblleft{}he was making\textquotedblright{}, the imperfect of \emph{facere}, \textquotedblleft{}to make, to do\textquotedblright{} — inherited, with the \emph{-ēba-} worn to \emph{-eva}, the same \emph{b} to \emph{v} as \emph{habēbat} and \textbf{aveva}. That \emph{facere} is the root of English \textbf{fact}, \textbf{factory} and \textbf{effect}. (Lewis \& Short, under \emph{facio}.)
\end{cousinweb}

\subsection*{Your turn}
Say these aloud:

\begin{itemize}
\raggedright
  \item \emph{Faceva freddo.} — then \emph{Faceva caldo.}
  \item \emph{pioveva} — \emph{piovere} on the plain recipe
  \item a thing, then a day — \emph{Il caffè era freddo.} then \emph{Faceva freddo.}
  \item a whole morning — \emph{Faceva freddo e pioveva. Non c'era il sole.}
\end{itemize}

\begin{tcolorbox}[breakable,colback=teal!4,colframe=teal!35!black,title={Before you move on}]
Say \textquotedblleft{}it was hot\textquotedblright{}. (\emph{\textbf{Faceva caldo.}}) Which do you use for the coffee, \emph{era} or \emph{faceva}? (\emph{\textbf{Era}} — it is about a thing.) Which Latin verb is inside \emph{faceva}, and what English word shares it? (\emph{\textbf{Facere}}, \textquotedblleft{}to make\textquotedblright{} — \textbf{fact}, \textbf{factory}.)
\end{tcolorbox}
\section[dormivo / ho dormito]{dormivo / ho dormito — the same sleep, told two ways}

\label{lesson:IT-C263-dormivo}



\begin{quote}
Paint a cold morning in two short sentences. (\emph{Faceva freddo. Non c'era il sole.})

Nothing has happened in it yet. This lesson is about what happens.
\end{quote}

\begin{grammarlens}[title={Grammar Lens: the scene and the events}]
Take one fact, that you slept, and let the tense decide what kind of fact it is:

\begin{quote}
\textbf{Dormivo.} — I was sleeping. \emph{(no edge: what was going on)} \textbf{Ho dormito.} — I slept. \emph{(an edge: a thing that happened and ended)}
\end{quote}

Neither is more correct. They answer different questions, and a told story needs both, in a division of labour:

\begin{quote}
The \textbf{imperfetto} paints the \textbf{scene}: what was going on, what things were like. The \textbf{passato prossimo} drops the \textbf{events} into it: what happened.
\end{quote}

\emph{Faceva freddo. Dormivo. Mia madre ha telefonato.} — It was cold. I was asleep. My mother phoned.

\textbf{Mentre}, \textquotedblleft{}while\textquotedblright{}, puts both into one sentence. \emph{Mentre dormivo, mia madre ha telefonato.} The \emph{mentre} half is the scene, still running; the other half is the event that cut into it.

English has no separate ending for this. It leans on \emph{was ‑ing} and \emph{used to} and lets the context carry the rest. Italian settles it in the ending.
\end{grammarlens}

\begin{cousinweb}
\textbf{imperfetto} is a learned word from Latin \emph{imperfectus}, \textquotedblleft{}unfinished\textquotedblright{} (EtymDB 2.1): \emph{in-}, \textquotedblleft{}not\textquotedblright{}, plus \emph{perfectus}, \textquotedblleft{}carried all the way through\textquotedblright{}, which is itself \emph{per-}, \textquotedblleft{}through\textquotedblright{}, plus \emph{facere}, \textquotedblleft{}to do, to make\textquotedblright{} (Lewis \& Short). That \emph{facere} is the one inside last lesson's \emph{faceva}.

So the name is a description of the job. The imperfetto is the \textbf{not-carried-through} past: unfinished in time, with no edge on it. English \textbf{perfect} still means finished all the way through.
\end{cousinweb}

\subsection*{Your turn}
Say these aloud:

\begin{itemize}
\raggedright
  \item the pair — \emph{dormivo}, then \emph{ho dormito}
  \item a scene, then an event — \emph{Pioveva. Ho aspettato l'autobus.}
  \item \emph{Mentre dormivo, mia madre ha telefonato.}
  \item \emph{C'era il sole, quindi ho camminato.}
\end{itemize}

\begin{tcolorbox}[breakable,colback=teal!4,colframe=teal!35!black,title={Before you move on}]
In a story, which past paints and which drops the events in? (The \emph{\textbf{imperfetto}} paints; the \emph{\textbf{passato prossimo}} drops them in.) What does \emph{imperfectus} say? (\textbf{Not carried through.}) In \emph{mentre dormivo, mia madre ha telefonato}, which half is the scene? (\emph{\textbf{Mentre dormivo.}})
\end{tcolorbox}
\section[(review)]{Review — the scene, and what cut into it}

\label{lesson:IT-C263-ripasso-mentre}



\begin{quote}
Build \textquotedblleft{}I was reading\textquotedblright{} from \emph{leggere}. (\emph{Leggevo.}) Then say \textquotedblleft{}I was tired\textquotedblright{}. (\emph{Ero stanco}, or \emph{ero stanca}.)
\end{quote}

\begin{grammarlens}[title={What you've built}]
Four lessons, and every piece of a scene is in your hands. The recipe, \textbf{-vo}, on any regular verb: \emph{parlavo}, \emph{leggevo}, \emph{dormivo}. The one verb that keeps its Latin shape: \emph{ero}, \emph{era}, \emph{c'era}. The weather, made rather than being: \emph{faceva freddo}. And the division of labour: the imperfetto paints, the passato prossimo drops the events in.

The joint between the two is \textbf{mentre}. Whatever follows it is the scene, still running; the rest of the sentence is the event.

\begin{quote}
\emph{Mentre leggevo, mia madre ha telefonato.} \emph{Mentre aspettavo l'autobus, ha cominciato a piovere.}
\end{quote}

The second one is worth hearing twice. The rain is an event here, because it \textbf{began}: \emph{ha cominciato}. Once it is falling, it is scene again: \emph{pioveva}.
\end{grammarlens}

\subsection*{Your turn}
Say these aloud:

\begin{itemize}
\raggedright
  \item a scene in three — \emph{Faceva freddo. C'era il sole. Ero stanco.}
  \item \emph{Mentre leggevo, mia madre ha telefonato.}
  \item the rain as an event — \emph{ha cominciato a piovere} — then as a scene — \emph{pioveva}
  \item one of your own — \emph{Mentre} plus an imperfetto, then one thing that happened
\end{itemize}

\begin{tcolorbox}[breakable,colback=teal!4,colframe=teal!35!black,title={Before you move on}]
After \emph{mentre}, which past? (The \textbf{imperfetto} — the scene.) Why is \emph{ha cominciato a piovere} in the passato prossimo? (The rain \textbf{began}: that is an event.) What does \emph{imperfetto} mean, taken apart? (\textbf{Not carried through.}) And where did its \emph{-v-} come from? (The \emph{\textbf{b}} of Latin \emph{amābam}.)
\end{tcolorbox}
